% TOP_PROBLEM: {"id": 30006871, "title": "Alperins Weight Conjecture for Modular Representations of Finite Groups", "category_id": 4, "category": {"id": 4, "name": "algebra", "display_name": "Algebra", "description": "Group theory, ring theory, field theory, and algebraic structures.", "slug": "algebra", "order_index": 4, "created_at": "2026-07-31T15:26:25.671Z"}, "status": "open"}
% FIRST_POSTED: 2026-10-06
% ATTEMPT_STATUS: solved
% OPENAI_CLAIM: {"schema_version":1,"source_repo":"https://github.com/openai/math","source_commit":"adc7f1241b42e322a6451854ab7e4b4c146bf78a","release_date":"2026-10-06","match":"full","resolution":"proved","independently_reviewed":false,"second_pass":"agreed","families":[{"family":"202","title":"The blockwise Alperin weight conjecture","subject":"Algebra","manuscripts":[{"title":"The Blockwise Alperin Weight Conjecture","pdf_path":"preprints/The-Blockwise-Alperin-Weight-Conjecture-September-23-2026/paper.pdf","readme_path":"preprints/The-Blockwise-Alperin-Weight-Conjecture-September-23-2026/README.md","date":"2026-09-23","theorem_ref":"Theorem 1.1"}],"lean":null,"reasoning_trace":null}]}
\documentclass[11pt]{article}
\usepackage[margin=25mm]{geometry}
\usepackage{fontspec}
\setmainfont{Latin Modern Roman}
\usepackage{amsmath,amssymb,mathtools}
\usepackage{unicode-math}
\setmathfont{Latin Modern Math}
\usepackage{enumitem}
\usepackage{xurl,hyperref,bookmark}
\hypersetup{colorlinks=true,urlcolor=blue}
\setcounter{secnumdepth}{0}
\setlength{\parindent}{0pt}
\setlength{\parskip}{5pt}
\setlength{\emergencystretch}{3em}
\begin{document}
\section*{Alperins Weight Conjecture for Modular Representations of Finite Groups}

\subsection{OpenAI's claim}
OpenAI presents a proof in the selected manuscript. For each prime and each block of a finite group, the claimed count of simple modular characters equals the count of block weights up to conjugacy.

\subsection{Outline of the argument}
The proof organizes alternating cancellations along subgroup chains and detects defect-zero contributions by central traces. A marked genus-one covering construction, valued lifting, and descent provide the uniform arithmetic input needed for the final chain-orbit cancellation.

\subsection{Scope relative to this problem}
The asserted equality counts Brauer characters and weights within each modular block. Its weights use the defect-zero normalizer quotient and central-character induction convention of this question, so the blockwise numerical assertion has the required scope. This is presented as a claimed affirmative answer.

\subsection{Formal verification}
Family 202 has no Lean scope page in this release. Its selected result is not presented as formalized. This site has not executed the Lean code. The scope descriptions do not certify the remaining manuscript arguments.

\subsection{Status caveat}
OpenAI attributes these manuscripts to work by an internal model. Its release cautions that arguments without formalization can contain mistakes. This catalogue has not independently checked the proofs; the labels report the mathematical claims attributed to OpenAI.

\subsection{Sources}
\begin{itemize}
\item OpenAI, \emph{The Blockwise Alperin Weight Conjecture}, 2026-09-23; citation key \texttt{OAI:The-Blockwise-Alperin-Weight-Conjecture-September-23-2026}. \href{https://github.com/openai/math/blob/adc7f1241b42e322a6451854ab7e4b4c146bf78a/preprints/The-Blockwise-Alperin-Weight-Conjecture-September-23-2026/paper.pdf}{PDF (GitHub)}; \href{https://github.com/openai/math/raw/adc7f1241b42e322a6451854ab7e4b4c146bf78a/preprints/The-Blockwise-Alperin-Weight-Conjecture-September-23-2026/paper.pdf}{PDF (direct)}; \href{https://github.com/openai/math/blob/adc7f1241b42e322a6451854ab7e4b4c146bf78a/preprints/The-Blockwise-Alperin-Weight-Conjecture-September-23-2026/README.md}{Manuscript page and citation}; \href{https://github.com/openai/math/blob/main/preprints/The-Blockwise-Alperin-Weight-Conjecture-September-23-2026/paper.pdf}{Latest version}.
\item \href{https://github.com/openai/math/blob/adc7f1241b42e322a6451854ab7e4b4c146bf78a/CONTENTS.md}{OpenAI catalogue entry}.
\end{itemize}
\end{document}
